
The APPS dataset contains reference solutions for 85.32\% of competition-level problems. A DeepSeek-Coder-7B model trained on those solutions with supervised fine-tuning achieves 0.0 pass@1 on LiveCodeBench-Hard. The training data exists; the generalization does not. This observation revises the standard mechanistic account of why RLEF outperforms SFT at hard benchmarks: the relevant void is a generalization void, not a dataset void.

Three contributions follow from this investigation. First, a controlled, reproducible, open-source pipeline for comparing RLEF against SFT across the full benchmark difficulty spectrum, using public components throughout. Second, empirical support for a statistically significant positive-ordered difficulty-scaling advantage for RLEF over SFT (JT z=+56.10, $p\approx 0)$, with the largest gains at LiveCodeBench-Hard ($\Delta =+0.18$, smoke-scale proxy), all values pending full-scale confirmation. Third, a controlled null result on reward formulation: fraction-of-tests and binary RLEF rewards produce equivalent performance at APPS training scale (p=0.552), providing a third independent replication of reward formulation convergence under GRPO on a different model and dataset combination.

The generalization void finding has implications beyond this study. The common assumption that harder benchmark performance improves by adding more hard training examples may not hold if the model cannot generalize those examples to held-out problems. Understanding this generalization gap, and confirming whether RLEF's self-generated feedback mechanism addresses it, is a central open question for the field.

